\textbf{Phần 1. Trắc nghiệm lựa chọn}
\setcounterpageref{ex}{0}

\begin{ex} %[5H9V51]
Một chiếc hộp không nắp hình lập phương có diện tích toàn phần là $180\text{ cm}^2$. Thể tích của chiếc hộp đó là
\choice
{$125\text{ cm}^3$}
{\True $216\text{ cm}^3$}
{$343\text{ cm}^3$}
{$64\text{ cm}^3$}
\loigiai{
Vì hộp không có nắp nên diện tích toàn phần của hộp gồm $5$ mặt vuông bằng nhau.\\
Diện tích của một mặt là:
$$180 : 5 = 36\text{ (cm}^2\text{)}.$$
Vì $6 \times 6 = 36$ nên cạnh của hình lập phương là $6\text{ cm}$.\\
Thể tích của chiếc hộp là:
$$6 \times 6 \times 6 = 216\text{ (cm}^3\text{)}.$$
\textbf{Đáp án của câu hỏi là B}
}
\end{ex}

\begin{ex} %[5H9V51]
Nếu cạnh của một hình lập phương gấp lên $3$ lần thì diện tích toàn phần của nó gấp lên bao nhiêu lần?
\choice
{$3\text{ lần}$}
{\True $9\text{ lần}$}
{$6\text{ lần}$}
{$27\text{ lần}$}
\loigiai{
Diện tích toàn phần ban đầu: $S_{tp} = a \times a \times 6$.\\
Khi cạnh gấp lên $3$ lần:
$$S'_{tp} = (3 \times a) \times (3 \times a) \times 6 = 9 \times (a \times a \times 6) = 9 \times S_{tp}.$$
Vậy diện tích toàn phần gấp lên $9$ lần.\\
\textbf{Đáp án của câu hỏi là B}
}
\end{ex}

\begin{ex} %[5H9V51]
Một hình lập phương có tổng diện tích xung quanh và diện tích toàn phần bằng $250\text{ cm}^2$. Cạnh của hình lập phương đó là
\choice
{$4\text{ cm}$}
{\True $5\text{ cm}$}
{$6\text{ cm}$}
{$25\text{ cm}$}
\loigiai{
Gọi diện tích một mặt là $S_1$.\\
Diện tích xung quanh là $4 \times S_1$, diện tích toàn phần là $6 \times S_1$.\\
Tổng diện tích xung quanh và toàn phần là:
$$4 \times S_1 + 6 \times S_1 = 10 \times S_1 = 250\text{ (cm}^2\text{)}.$$
Diện tích một mặt là:
$$S_1 = 250 : 10 = 25\text{ (cm}^2\text{)}.$$
Vì $5 \times 5 = 25$ nên độ dài cạnh hình lập phương là $5\text{ cm}$.\\
\textbf{Đáp án của câu hỏi là B}
}
\end{ex}

\begin{ex} %[5H9V51]
Người ta làm một khung lồng chim dạng hình lập phương bằng sắt, mỗi cạnh dài $40\text{ cm}$. Người ta bọc lưới thép quanh $4$ mặt bên của lồng chim. Diện tích lưới thép cần dùng là
\choice
{$160\text{ dm}^2$}
{\True $64\text{ dm}^2$}
{$96\text{ dm}^2$}
{$640\text{ dm}^2$}
\loigiai{
Đổi: $40\text{ cm} = 4\text{ dm}$.\\
Bọc lưới thép quanh $4$ mặt bên chính là diện tích xung quanh của hình lập phương:\\
$$S_{xq} = 4 \times 4 \times 4 = 64\text{ (dm}^2\text{)}.$$
\textbf{Đáp án của câu hỏi là B}
}
\end{ex}

\begin{ex} %[5H9V51]
Một cái thùng hình lập phương có cạnh $1{,}2\text{ m}$. Người ta sơn toàn bộ mặt trong và mặt ngoài của chiếc thùng đó (thùng có nắp kín). Biết giá tiền sơn mỗi mét vuông là $30\,000\text{ đồng}$. Tổng số tiền sơn thùng là
\choice
{$259\,200\text{ đồng}$}
{\True $518\,400\text{ đồng}$}
{$345\,600\text{ đồng}$}
{$432\,000\text{ đồng}$}
\loigiai{
Diện tích toàn phần một mặt ngoài của thùng là:
$$1{,}2 \times 1{,}2 \times 6 = 8{,}64\text{ (m}^2\text{)}.$$
Tổng diện tích quét sơn cả mặt trong và mặt ngoài là:
$$8{,}64 \times 2 = 17{,}28\text{ (m}^2\text{)}.$$
Số tiền sơn thùng là:
$$17{,}28 \times 30\,000 = 518\,400\text{ (đồng)}.$$
\textbf{Đáp án của câu hỏi là B}
}
\end{ex}

\begin{ex} %[5H9V51]
Biết diện tích toàn phần của hình lập phương thứ nhất bằng $150\text{ dm}^2$, diện tích toàn phần của hình lập phương thứ hai bằng $600\text{ dm}^2$. Cạnh của hình lập phương thứ hai dài gấp mấy lần cạnh của hình lập phương thứ nhất?
\choice
{$4\text{ lần}$}
{\True $2\text{ lần}$}
{$3\text{ lần}$}
{$8\text{ lần}$}
\loigiai{
Diện tích một mặt hình lập phương thứ nhất là:
$$150 : 6 = 25\text{ (dm}^2\text{)} \Rightarrow \text{cạnh là } 5\text{ dm}.$$
Diện tích một mặt hình lập phương thứ hai là:
$$600 : 6 = 100\text{ (dm}^2\text{)} \Rightarrow \text{cạnh là } 10\text{ dm}.$$
Cạnh hình thứ hai gấp cạnh hình thứ nhất số lần là:
$$10 : 5 = 2\text{ (lần)}.$$
\textbf{Đáp án của câu hỏi là B}
}
\end{ex}

\begin{ex} %[5H9V51]
Ghép $8$ khối lập phương nhỏ cạnh $1\text{ cm}$ thành một khối lập phương lớn. Diện tích toàn phần của khối lập phương lớn là
\choice
{$48\text{ cm}^2$}
{\True $24\text{ cm}^2$}
{$36\text{ cm}^2$}
{$16\text{ cm}^2$}
\loigiai{
Khối lập phương lớn gồm $2 \times 2 \times 2 = 8$ khối nhỏ, nên cạnh của khối lập phương lớn là:
$$1 \times 2 = 2\text{ (cm)}.$$
Diện tích toàn phần của khối lập phương lớn là:
$$2 \times 2 \times 6 = 24\text{ (cm}^2\text{)}.$$
\textbf{Đáp án của câu hỏi là B}
}
\end{ex}

\begin{ex} %[5H9V51]
Một khối gỗ hình lập phương có diện tích toàn phần lớn hơn diện tích xung quanh là $72\text{ cm}^2$. Cạnh của khối gỗ đó là
\choice
{$36\text{ cm}$}
{\True $6\text{ cm}$}
{$12\text{ cm}$}
{$9\text{ cm}$}
\loigiai{
Diện tích toàn phần lớn hơn diện tích xung quanh chính là tổng diện tích hai mặt đáy của hình lập phương.\\
Diện tích của hai mặt đáy là $72\text{ cm}^2$.\\
Diện tích một mặt của hình lập phương là:
$$72 : 2 = 36\text{ (cm}^2\text{)}.$$
Vì $6 \times 6 = 36$ nên cạnh của khối gỗ là $6\text{ cm}$.\\
\textbf{Đáp án của câu hỏi là B}
}
\end{ex}

\begin{ex} %[5H9V51]
Người ta xếp các khối lập phương nhỏ cạnh $1\text{ cm}$ thành một khối lập phương lớn có cạnh $4\text{ cm}$, sau đó sơn tất cả các mặt ngoài của khối lớn. Hỏi có bao nhiêu khối lập phương nhỏ được sơn đúng $2$ mặt?
\choice
{$16$}
{\True $24$}
{$8$}
{$32$}
\loigiai{
Khối lập phương có cạnh $4\text{ cm}$ được ghép từ $4$ khối nhỏ mỗi cạnh.\\
Các khối lập phương nhỏ được sơn $2$ mặt nằm ở các cạnh của khối lớn (không kể các khối ở $8$ đỉnh).\\
Mỗi cạnh có số khối được sơn $2$ mặt là:
$$4 - 2 = 2\text{ (khối)}.$$
Hình lập phương có $12$ cạnh, nên tổng số khối được sơn đúng $2$ mặt là:
$$2 \times 12 = 24\text{ (khối)}.$$
\textbf{Đáp án của câu hỏi là B}
}
\end{ex}

\begin{ex} %[5H9V51]
Một chiếc hộp hình lập phương cạnh $10\text{ cm}$. Bạn Nam dán giấy màu đỏ vào $4$ mặt xung quanh và giấy màu xanh vào $2$ mặt đáy. Diện tích giấy màu đỏ nhiều hơn diện tích giấy màu xanh là
\choice
{$100\text{ cm}^2$}
{\True $200\text{ cm}^2$}
{$300\text{ cm}^2$}
{$400\text{ cm}^2$}
\loigiai{
Diện tích một mặt là:
$$10 \times 10 = 100\text{ (cm}^2\text{)}.$$
Diện tích giấy màu đỏ ($4$ mặt xung quanh) là:
$$100 \times 4 = 400\text{ (cm}^2\text{)}.$$
Diện tích giấy màu xanh ($2$ mặt đáy) là:
$$100 \times 2 = 200\text{ (cm}^2\text{)}.$$
Diện tích giấy đỏ nhiều hơn giấy xanh là:
$$400 - 200 = 200\text{ (cm}^2\text{)}.$$
\textbf{Đáp án của câu hỏi là B}
}
\end{ex}

\begin{ex} %[5H9V51]
Cho hai hình lập phương. Hình thứ nhất có cạnh bằng $2\text{ dm}$, hình thứ hai có cạnh bằng $\dfrac{1}{2}$ cạnh hình thứ nhất. Tỉ số giữa diện tích xung quanh của hình thứ hai và diện tích xung quanh của hình thứ nhất là
\choice
{$\dfrac{1}{2}$}
{\True $\dfrac{1}{4}$}
{$\dfrac{1}{8}$}
{$\dfrac{1}{16}$}
\loigiai{
Tỉ số giữa cạnh hình thứ hai và cạnh hình thứ nhất là $\dfrac{1}{2}$.\\
Tỉ số diện tích xung quanh giữa hình thứ hai và hình thứ nhất bằng bình phương tỉ số hai cạnh:\\
$$\left(\dfrac{1}{2}\right) \times \left(\dfrac{1}{2}\right) = \dfrac{1}{4}.$$
\textbf{Đáp án của câu hỏi là B}
}
\end{ex}

\begin{ex} %[5H9V51]
Người ta cắt một khối gỗ hình lập phương có diện tích toàn phần $96\text{ dm}^2$ thành $8$ khối lập phương nhỏ bằng nhau. Tổng diện tích toàn phần của cả $8$ khối lập phương nhỏ đó là
\choice
{$96\text{ dm}^2$}
{\True $192\text{ dm}^2$}
{$144\text{ dm}^2$}
{$288\text{ dm}^2$}
\loigiai{
Diện tích một mặt của khối gỗ ban đầu là:
$$96 : 6 = 16\text{ (dm}^2\text{)} \Rightarrow \text{cạnh là } 4\text{ dm}.$$
Cắt thành $8$ khối nhỏ bằng nhau thì mỗi khối nhỏ có cạnh là:
$$4 : 2 = 2\text{ (dm)}.$$
Diện tích toàn phần của một khối nhỏ là:
$$2 \times 2 \times 6 = 24\text{ (dm}^2\text{)}.$$
Tổng diện tích toàn phần của cả $8$ khối nhỏ là:
$$24 \times 8 = 192\text{ (dm}^2\text{)}.$$
\textbf{Đáp án của câu hỏi là B}
}
\end{ex}

\textbf{Phần 2. Trắc nghiệm đúng sai}
\setcounterpageref{ex}{0}

\begin{ex} %[5H9V51]
Cho một hình lập phương có độ dài cạnh là $a = 8\text{ cm}$. Xét tính đúng/sai của các nhận định sau:
\choiceTF
{\True Diện tích một mặt của hình lập phương là $64\text{ cm}^2$}
{\True Diện tích xung quanh của hình lập phương là $256\text{ cm}^2$}
{Diện tích toàn phần của hình lập phương là $512\text{ cm}^2$}
{\True Tỉ số giữa diện tích xung quanh và diện tích toàn phần của hình lập phương bằng $\dfrac{2}{3}$}
\loigiai{
\begin{itemchoice}
\item Đúng. $S_1 = 8 \times 8 = 64\text{ (cm}^2\text{)}$.
\item Đúng. $S_{xq} = 64 \times 4 = 256\text{ (cm}^2\text{)}$.
\item Sai. Diện tích toàn phần là $S_{tp} = 64 \times 6 = 384\text{ (cm}^2\text{)}$ (không phải $512\text{ cm}^2$).
\item Đúng. Tỉ số $\dfrac{S_{xq}}{S_{tp}} = \dfrac{4 \times S_1}{6 \times S_1} = \dfrac{4}{6} = \dfrac{2}{3}$.
\end{itemchoice}
}
\end{ex}

\begin{ex} %[5H9V51]
Một khối gỗ hình lập phương cạnh $3\text{ dm}$ được sơn kín tất cả các mặt ngoài rồi cắt thành $27$ khối lập phương nhỏ cạnh $1\text{ dm}$. Xét tính đúng/sai của các nhận định sau:
\choiceTF
{\True Có đúng $8$ khối lập phương nhỏ được sơn $3$ mặt}
{\True Có đúng $12$ khối lập phương nhỏ được sơn $2$ mặt}
{Có $8$ khối lập phương nhỏ được sơn đúng $1$ mặt}
{\True Có đúng $1$ khối lập phương nhỏ không được sơn mặt nào}
\loigiai{
\begin{itemchoice}
\item Đúng. $8$ khối nhỏ ở $8$ đỉnh được sơn đúng $3$ mặt.
\item Đúng. Số khối nhỏ ở các cạnh được sơn $2$ mặt là $(3 - 2) \times 12 = 12$ (khối).
\item Sai. Số khối nhỏ ở tâm các mặt được sơn $1$ mặt là $(3 - 2) \times (3 - 2) \times 6 = 6$ (khối), không phải $8$ khối.
\item Đúng. Số khối nhỏ ở lõi bên trong không được sơn mặt nào là $(3 - 2)^3 = 1$ (khối).
\end{itemchoice}
}
\end{ex}

\textbf{Phần 3. Câu trả lời ngắn}
\setcounterpageref{ex}{0}

\begin{ex} %[5H9V51]
Một cái hộp hình lập phương không nắp có diện tích toàn phần bằng $500\text{ cm}^2$. Diện tích xung quanh của chiếc hộp đó là bao nhiêu xăng-ti-mét vuông?
\shortans{400}
\loigiai{
Vì hộp không nắp nên diện tích toàn phần bằng $5$ lần diện tích một mặt.\\
Diện tích một mặt của hình lập phương là:
$$500 : 5 = 100\text{ (cm}^2\text{)}.$$
Diện tích xung quanh của chiếc hộp là:
$$100 \times 4 = 400\text{ (cm}^2\text{)}.$$
}
\end{ex}

\begin{ex} %[5H9V51]
Một hình lập phương có cạnh $7\text{ cm}$. Diện tích toàn phần của hình lập phương đó là bao nhiêu xăng-ti-mét vuông?
\shortans{294}
\loigiai{
Diện tích toàn phần của hình lập phương là:
$$7 \times 7 \times 6 = 294\text{ (cm}^2\text{)}.$$
}
\end{ex}

\begin{ex} %[5H9V51]
Biết diện tích xung quanh của một hình lập phương là $144\text{ dm}^2$. Cạnh của hình lập phương đó là bao nhiêu đề-xi-mét?
\shortans{6}
\loigiai{
Diện tích một mặt của hình lập phương là:
$$144 : 4 = 36\text{ (dm}^2\text{)}.$$
Vì $6 \times 6 = 36$ nên cạnh của hình lập phương là $6\text{ dm}$.
}
\end{ex}

\begin{ex} %[5H9V51]
Người ta ghép $4$ khối lập phương nhỏ cạnh $2\text{ cm}$ thành một khối hình chữ L trên một mặt phẳng. Tổng diện tích toàn phần của khối mới là bao nhiêu xăng-ti-mét vuông?
\shortans{72}
\loigiai{
Tổng số mặt ngoài của $4$ khối lập phương rời rạc là: $4 \times 6 = 24$ (mặt vuông).\\
Để ghép liền $4$ khối thành hình chữ L, cần $3$ mối ghép nối tiếp. Mỗi mối ghép nối che đi $2$ mặt vuông, do đó che mất:
$$3 \times 2 = 6\text{ (mặt)}.$$
Số mặt ngoài lộ ra của khối mới là:
$$24 - 6 = 18\text{ (mặt)}.$$
Diện tích mỗi mặt vuông là:
$$2 \times 2 = 4\text{ (cm}^2\text{)}.$$
Diện tích toàn phần của khối hình mới là:
$$18 \times 4 = 72\text{ (cm}^2\text{)}.$$
}
\end{ex}

\begin{ex} %[5H9V51]
Bác An làm một chiếc thùng tôn hình lập phương có nắp cạnh $1{,}5\text{ m}$. Cứ mỗi mét vuông tôn có giá $90\,000\text{ đồng}$. Hỏi bác An tốn bao nhiêu nghìn đồng để mua tôn (coi mép gò không đáng kể)?
\shortans{1215}
\loigiai{
Diện tích toàn phần của chiếc thùng tôn hình lập phương có nắp là:
$$1{,}5 \times 1{,}5 \times 6 = 13{,}5\text{ (m}^2\text{)}.$$
Số tiền mua tôn là:
$$13{,}5 \times 90\,000 = 1\,215\,000\text{ (đồng)} = 1\,215\text{ nghìn đồng}.$$
}
\end{ex}

\begin{ex} %[5H9V51]
Một hình lập phương có chu vi của một mặt đáy bằng $36\text{ cm}$. Thể tích của hình lập phương đó là bao nhiêu xăng-ti-mét khối?
\shortans{729}
\loigiai{
Cạnh của hình lập phương là:
$$36 : 4 = 9\text{ (cm)}.$$
Thể tích của hình lập phương là:
$$9 \times 9 \times 9 = 729\text{ (cm}^3\text{)}.$$
}
\end{ex}
